% Positive-test posterior with explicit prevalence.
\clearpage
\section*{Verification case: likelihood is not posterior probability}
\noindent Suppose a test detects a condition with sensitivity
\(P(+\mid D)=0.95\) and false-positive probability
\(P(+\mid\neg D)=0.05\). Assume the condition occurs in
\(P(D)=0.01\) of the population. Then Bayes' rule gives
\[
P(D\mid+)
=\frac{P(+\mid D)P(D)}
{P(+\mid D)P(D)+P(+\mid\neg D)P(\neg D)}
=\frac{0.0095}{0.0095+0.0495}
\approx0.1610.
\]
A positive test means about \(16.1\%\) posterior probability
under \emph{this declared model}, not \(95\%\).
The contrast is due to prevalence, not an algebraic inconsistency.

For a hypothetical population of \(10{,}000\), the model expects
\(100\) affected people, \(95\) true positives and \(495\)
false positives among the \(9{,}900\) unaffected.
Thus \(95/(95+495)=95/590\approx0.1610\).
These are expected counts from the model, not a real dataset.

\paragraph{Use in visual computing.}
A pose-estimation failure detector or rendering anomaly classifier
has the same dependence on the prior frequency of failures.
Changing the deployment population changes predictive values
even if sensitivity and false-positive rates remain fixed.
Classifiers must therefore report prevalence, calibration and
the distinction between conditional likelihood and posterior
confidence.

\paragraph{Reconstruction.}
Implement Bayes' rule and compare the direct algebra with
expected-count reasoning. Sweep prevalence over logarithmically
spaced values and plot the posterior. Test both impossible
evidence (\(P(+)=0\)) and inconsistent input probabilities as
explicit invalid cases; no undefined division should produce a
plausible confidence score.
